\documentclass[a4paper,11pt]{jsarticle}
\usepackage{geometry}
\geometry{top=22mm, bottom=24mm, left=22mm, right=22mm}
\usepackage{longtable}
\usepackage{array}
\setlength{\arrayrulewidth}{0.6pt}
\renewcommand{\arraystretch}{1.18}
\usepackage{listings}
\usepackage[dvipdfmx]{xcolor}
\usepackage[dvipdfmx]{hyperref}
\hypersetup{colorlinks=true, linkcolor=black, urlcolor=blue}

\lstdefinestyle{code}{
  language=Java,
  basicstyle=\ttfamily\small,
  keywordstyle=\color{blue!70!black}\bfseries,
  commentstyle=\color{gray!70!black},
  stringstyle=\color{red!60!black},
  numbers=left,
  numberstyle=\tiny\color{gray},
  breaklines=true,
  frame=single,
  tabsize=2,
  captionpos=b,
  showstringspaces=false
}

\lstdefinestyle{textdiagram}{
  basicstyle=\ttfamily\small,
  numbers=none,
  breaklines=true,
  frame=single,
  tabsize=2,
  showstringspaces=false
}

\title{チーム23 ハッカソン1\\
Processing 音源サブシステム\\
基本設計・詳細設計 仕様書}
\author{梅澤颯太（学籍番号：25G1021）}
\date{2026年4月29日 改訂案}

\begin{document}
\maketitle
\tableofcontents
\newpage

\section{基本設計：機能一覧とシステム構成図}

\subsection{機能一覧}

Processing音源サブシステムの機能一覧を表\ref{tab:fbs}に示す．
塩澤設計のF4「発音情報配信」に対して，本書ではPC側の受信境界から音声出力までを
F6として定義する．

\begin{longtable}{p{0.28\linewidth}p{0.62\linewidth}}
\caption{Processing音源サブシステムの機能一覧}\label{tab:fbs}\\
\hline
機能 & 役割 \\
\hline
\endfirsthead
\hline
機能 & 役割 \\
\hline
\endhead
Serial接続管理 &
Processing起動時に1つのUSB Serialポートを開く．切断時は再接続待ちに遷移し，
ユーザーがポートを再選択できるようにする． \\
\hline
NOTEフレーム同期・復号 &
Serialから届くバイト列を内部バッファへ蓄積し，magic=0x4F52を探して20B固定長の
NOTEパケットへ復号する．version/type/seq/partIdを検査する． \\
\hline
パート設定管理 &
Mac 4台構成では各Processingが自パート1つだけを鳴らす．統合デモや検証用として，
1台で4パートを受けるモードも持つ． \\
\hline
ボイス管理 &
NoteOn/NoteOff，durationMs経過による自動停止，同一partId内の重複発音，
同時発音数上限を管理する． \\
\hline
音色生成 &
partIdに応じて金管波形またはドラム波形を選択し，ADSRエンベロープを適用して
MinimのAudioOutputへ送る． \\
\hline
画面表示・操作IF &
接続状態，partId，seq，受信間隔，発音中ボイス数，音量，ミュート状態を表示し，
キー操作でポート再接続，ミュート，テスト発音を行う． \\
\hline
ログ・検証 &
受信時刻，発音時刻，seq欠落，破棄理由をCSVログとして保存し，結合試験で
遅延・取りこぼし・音色確認に使えるようにする． \\
\hline
\end{longtable}

\subsection{ユーザー目線のシステム構成}

本書の主構成は，塩澤設計に合わせて「楽器ノード1台につきMac 1台」である．
各Mac上のProcessingは，自分にUSB接続されたArduino UNO R4 WiFiからNOTEだけを受け，
自パートの音を鳴らす．

\begin{lstlisting}[style=textdiagram, caption={Processing音源サブシステムを含む全体構成}]
User conducts
     |
     v
node_01: Conductor, SoftAP, CTRL/BEAT sender
     |
     | UDP multicast CTRL/BEAT
     v
+---------+   USB Serial NOTE   +------------------+   Audio out
| node_02 | -------------------> | Mac 1 Processing | ----------> Brass 1
| node_03 | -------------------> | Mac 2 Processing | ----------> Brass 2
| node_04 | -------------------> | Mac 3 Processing | ----------> Brass 3
| node_05 | -------------------> | Mac 4 Processing | ----------> Drum
+---------+                      +------------------+
\end{lstlisting}

\subsection{パート対応}

\begin{table}[htbp]
\centering
\caption{partIdとProcessing音色の対応}
\label{tab:part}
\begin{tabular}{ccll}
\hline
ノード & partId & 役割 & Processing音色 \\
\hline
node\_02 & 0x02 & 金管1（主旋律） & BrassVoice \\
\hline
node\_03 & 0x03 & 金管2（輪唱1，4拍遅れ） & BrassVoice \\
\hline
node\_04 & 0x04 & 金管3（輪唱2，8拍遅れ） & BrassVoice \\
\hline
node\_05 & 0x05 & ドラム（リズム伴奏） & DrumVoice \\
\hline
\end{tabular}
\end{table}

Mac 4台構成では，partIdは接続確認用としても利用する．
たとえばMac 1は期待partIdを0x02に設定し，0x03--0x05を受信した場合は
画面に警告を表示して発音しない．
統合デモモードでは期待partIdを\texttt{ALL}とし，4パートすべてを鳴らす．

\subsection{性能指標}

\begin{longtable}{p{0.13\linewidth}p{0.36\linewidth}p{0.38\linewidth}}
\caption{Processing側MOP}\label{tab:mop}\\
\hline
ID & 指標 & 目標値 \\
\hline
\endfirsthead
\hline
ID & 指標 & 目標値 \\
\hline
\endhead
PMOP-1 & NOTE受信から発音予約までの処理遅延 & 5 ms以下 \\
\hline
PMOP-2 & Serialフレーム同期復帰 & 破損バイト混入後，次の正常magicから1パケット以内に復帰 \\
\hline
PMOP-3 & partId誤接続検出 & 期待partIdと異なるNOTEを100\%警告表示し，発音しない \\
\hline
PMOP-4 & 同時発音管理 & 金管は各パート最大4ボイス，ドラムは最大8ボイスまで破綻なく再生 \\
\hline
PMOP-5 & 音量安全性 & velocity=127でもAudioOutputがクリップしないよう総振幅を0.8以下に正規化 \\
\hline
PMOP-6 & 音色根拠 & 金管の倍音比率とADSR値を表として管理し，FFT解析結果で更新可能 \\
\hline
\end{longtable}

\section{基本設計：操作インターフェースと状態遷移}

\subsection{ユーザー操作}

Processing側は，デモ担当者または演奏準備担当者が操作する．
観客は直接PCを操作せず，指揮者の動きにより間接的にテンポと強弱を変える．

\begin{longtable}{p{0.22\linewidth}p{0.30\linewidth}p{0.36\linewidth}}
\hline
操作 & 対象 & システム挙動 \\
\hline
\endhead
Processing起動 & Mac担当者 & Serialポート一覧を表示し，前回保存したポートがあれば自動接続を試みる． \\
\hline
キー\texttt{1}--\texttt{4} & Mac担当者 & 期待partIdを0x02--0x05に切り替える．本番ではMacごとに固定する． \\
\hline
キー\texttt{a} & Mac担当者 & 統合デモモードに入り，partId 0x02--0x05をすべて受け付ける． \\
\hline
キー\texttt{r} & Mac担当者 & Serialを閉じ，ポート再接続状態へ戻る． \\
\hline
キー\texttt{m} & Mac担当者 & ミュートを切り替える．ミュート中も受信ログは残す． \\
\hline
キー\texttt{t} & Mac担当者 & 選択中partIdのテスト音を鳴らし，スピーカ接続と音量を確認する． \\
\hline
キー\texttt{l} & Mac担当者 & CSVログ保存を開始・停止する． \\
\hline
指揮動作 & 指揮者 & Arduino側がテンポ・velocityを変化させ，NOTEのdurationMsとvelocityに反映される． \\
\hline
\end{longtable}

\subsection{画面表示}

Processing画面は，演奏中の障害を即座に見つけるための最小UIとする．
表示項目は以下である．

\begin{itemize}
\item 接続状態：Disconnected / PortSelect / Ready / Playing / Error
\item 接続ポート名，期待partId，実際に受信したpartId
\item 直近seq，seq欠落数，破棄パケット数
\item 直近NOTEのnoteNumber，velocity，gate，durationMs
\item 発音中ボイス数，ミュート状態，ログ保存状態
\item 波形モニタ（既存の\texttt{pc\_app/processing\_test}の表示を流用）
\end{itemize}

\subsection{状態遷移}

\begin{lstlisting}[style=textdiagram, caption={Processing音源サブシステムの状態遷移}]
Boot
  -> PortSelect      : setup() completed, serial ports listed
PortSelect
  -> Ready           : selected port opened
  -> Error           : port open failed
Ready
  -> Playing         : first valid NOTE received
  -> PortSelect      : user presses r
Playing
  -> Muted           : user presses m
  -> Error           : serial exception or timeout > 3000 ms
  -> Ready           : no NOTE for 1000 ms, connection still alive
Muted
  -> Playing         : user presses m
  -> Error           : serial exception
Error
  -> PortSelect      : user presses r or auto retry interval elapsed
\end{lstlisting}

\begin{longtable}{p{0.18\linewidth}p{0.36\linewidth}p{0.34\linewidth}}
\hline
状態 & 意味 & 主な処理 \\
\hline
\endhead
Boot & Processing起動直後 & Minim，AudioOutput，設定値，ログファイルを初期化する． \\
\hline
PortSelect & Serialポート選択待ち & 利用可能ポートを画面表示し，設定済みポートへ接続を試みる． \\
\hline
Ready & 接続済み・NOTE待ち & Serialバッファを空にし，magic探索を開始する． \\
\hline
Playing & 正常受信・発音中 & NOTE復号，partId検査，Voice生成，画面更新，ログ出力を行う． \\
\hline
Muted & 受信継続・発音停止 & NOTEは復号しログに残すが，Voiceは生成しない． \\
\hline
Error & 接続異常・復号不能 & Serialを閉じ，発音中Voiceを停止し，破棄理由を表示する． \\
\hline
\end{longtable}

\section{詳細設計：ハードウェア設計}

Processing担当範囲の主対象はソフトウェアであるが，ルーブリックのハードウェア設計項目に対応するため，
PC側に接続される構成要素・接続関係・動作を明確化する．

\subsection{構成要素}

\begin{longtable}{p{0.18\linewidth}p{0.20\linewidth}p{0.10\linewidth}p{0.40\linewidth}}
\hline
区分 & 要素 & 数量 & 役割 \\
\hline
\endhead
H1 & Arduino UNO R4 WiFi & 4 & node\_02--05．楽譜進行後，NOTEパケットをUSB Serialで送信する． \\
\hline
H2 & Mac またはPC & 4 & 各楽器ノード専属のProcessing実行機．1台につき1Serialポートだけ開く． \\
\hline
H3 & USB Type-Cケーブル & 4 & 各Arduinoと各Macを1対1で接続し，給電とSerial通信を兼ねる． \\
\hline
H4 & スピーカまたは内蔵スピーカ & 4 & ProcessingのAudioOutputを再生する．本番では各Macの音量を同程度に調整する． \\
\hline
H5 & 指揮者ノード & 1 & node\_01．Processingには直接接続しないが，BEAT/CTRLによりNOTE発生時刻を決める． \\
\hline
\end{longtable}

\subsection{接続設計}

\begin{lstlisting}[style=textdiagram, caption={Mac 4台構成の接続図}]
node_02 USB-C ---- Mac 1 ---- speaker 1 : partId 0x02, Brass 1
node_03 USB-C ---- Mac 2 ---- speaker 2 : partId 0x03, Brass 2
node_04 USB-C ---- Mac 3 ---- speaker 3 : partId 0x04, Brass 3
node_05 USB-C ---- Mac 4 ---- speaker 4 : partId 0x05, Drum
\end{lstlisting}

\begin{longtable}{p{0.24\linewidth}p{0.64\linewidth}}
\hline
項目 & 設計 \\
\hline
\endhead
通信方式 & USB CDC Serial．Arduino側は\texttt{Serial.write()}で20B固定長バイナリを書き込む． \\
\hline
ボーレート & 115200 bps．USB CDCでは実効速度は十分だが，Arduino側の\texttt{NoteSenderConfig}と統一する． \\
\hline
フレーミング & 先頭2Bのmagic=0x4F52を同期マークとする．Serial上のバイト列はリトルエンディアンなので\texttt{0x52 0x4F}の順で現れる． \\
\hline
音声出力 & ProcessingのMinim \texttt{AudioOutput}を各Macの既定出力へ送る．会場ではOS音量を70\%程度から開始し，クリップがないよう調整する． \\
\hline
誤接続対策 & Macごとに期待partIdを固定し，異なるpartIdを受信したら画面表示・ログ出力し発音しない． \\
\hline
故障時動作 & USBが抜けた場合はErrorへ遷移し，すべての発音を停止する．再接続後，キー\texttt{r}でPortSelectへ戻す． \\
\hline
\end{longtable}

\section{詳細設計：ソフトウェア設計}

\subsection{ファイル構成}

実装時は，\texttt{pc\_app/orchestra\_processing/}にProcessingスケッチを置く．
Processingの仕様上，フォルダ名とメイン\texttt{.pde}名は一致させる．

\begin{lstlisting}[style=textdiagram, caption={Processing実装予定ファイル構成}]
pc_app/
`-- orchestra_processing/
    |-- orchestra_processing.pde   // setup(), draw(), keyPressed()
    |-- Config.pde                  // partId, serial port, audio params
    |-- OrcProtocol.pde             // packet structs, little-endian decoder
    |-- SerialFrameReader.pde       // byte buffer and magic sync
    |-- PartManager.pde             // part filter and voice allocation
    |-- Voice.pde                   // BrassVoice, DrumVoice, ADSR
    |-- ToneFactory.pde             // waveform tables
    |-- UiView.pde                  // status display and waveform monitor
    `-- Logger.pde                  // CSV logging
\end{lstlisting}

\subsection{NOTEパケット定義}

塩澤設計の\texttt{OrcProtocol}に合わせ，Processing側は表\ref{tab:note-packet}の20Bを読む．
CTRL/BEATはProcessingへ直接来ないため，type=1またはtype=2を受信した場合は破棄する．

\begin{longtable}{p{0.12\linewidth}p{0.20\linewidth}p{0.16\linewidth}p{0.40\linewidth}}
\caption{NOTEパケット}\label{tab:note-packet}\\
\hline
offset & field & type & Processing側の扱い \\
\hline
\endfirsthead
\hline
offset & field & type & Processing側の扱い \\
\hline
\endhead
0 & magic & uint16 & 0x4F52以外なら同期探索を続ける． \\
\hline
2 & version & uint8 & 0x01以外なら破棄し，version errorとしてログする． \\
\hline
3 & type & uint8 & 3のときのみNOTEとして処理する． \\
\hline
4 & seq & uint32 & 前回seqとの差を取り，欠落・重複を検出する． \\
\hline
8 & timestampMs & uint32 & Arduino側送信時刻．ログに保存し，遅延評価に使う． \\
\hline
12 & partId & uint8 & 0x02--0x05．期待partIdまたはALLと一致する場合だけ発音する． \\
\hline
13 & noteNumber & uint8 & MIDIノート番号．0は休符またはNoteOff相当として扱う． \\
\hline
14 & velocity & uint8 & 0--127．発音振幅へ変換する． \\
\hline
15 & gate & uint8 & 1=NoteOn，0=NoteOff．その他の値は破棄する． \\
\hline
16 & durationMs & uint16 & NoteOn時の自動NoteOff予定時刻を計算する．0なら最小長50 msを使う． \\
\hline
18 & reserved & uint8[2] & 0以外でも発音は継続するが，ログに警告を残す． \\
\hline
\end{longtable}

\subsection{クラス設計}

\begin{longtable}{p{0.21\linewidth}p{0.31\linewidth}p{0.36\linewidth}}
\hline
クラス/構造体 & 主なフィールド・メソッド & 役割 \\
\hline
\endhead
\texttt{NotePacket} &
\texttt{partId, noteNumber, velocity, gate, durationMs, seq, timestampMs} &
20Bバイナリから復号した発音指令を表す． \\
\hline
\texttt{SerialFrameReader} &
\texttt{pushByte()}, \texttt{tryReadPacket()}, \texttt{dropUntilMagic()} &
Serial入力を蓄積し，magic同期と20B切り出しを行う．破損時は1Bずつ捨てて再同期する． \\
\hline
\texttt{PartConfig} &
\texttt{expectedPartId}, \texttt{instrumentType}, \texttt{maxVoices} &
Macごとの担当パートと音色を定義する． \\
\hline
\texttt{PartManager} &
\texttt{handleNote()}, \texttt{noteOn()}, \texttt{noteOff()}, \texttt{update()} &
partId検査，同時発音上限，durationMs経過による停止を管理する． \\
\hline
\texttt{Voice} &
\texttt{start()}, \texttt{release()}, \texttt{isFinished()} &
1つの発音単位．Minimの\texttt{Oscil}とエンベロープを保持する． \\
\hline
\texttt{BrassVoice} &
\texttt{makeBrassWaveform()}, \texttt{attackMs}, \texttt{releaseMs} &
金管1--3の音色を作る．倍音テーブルとADSR値は表で管理する． \\
\hline
\texttt{DrumVoice} &
\texttt{makeNoiseBurst()}, \texttt{decayMs} &
node\_05用．ノイズ成分と短い減衰でリズム伴奏を表現する． \\
\hline
\texttt{UiView} &
\texttt{drawStatus()}, \texttt{drawWaveform()} &
状態，受信内容，エラー，波形モニタを描画する． \\
\hline
\texttt{Logger} &
\texttt{logPacket()}, \texttt{logError()}, \texttt{flush()} &
結合試験用CSVを保存する． \\
\hline
\end{longtable}

\subsection{音色パラメータ}

\subsubsection{金管音色}

金管音色は，加算合成による倍音テーブルとADSRで作る．
初版では\texttt{Line}による単純な直線減衰としていたが，
金管らしさを出すため，本設計では最初からADSR相当のエンベロープを採用する．

\begin{table}[htbp]
\centering
\caption{金管音色の初期倍音比率}
\begin{tabular}{cccccccc}
\hline
倍音 & 1 & 2 & 3 & 4 & 5 & 6 & 7 \\
\hline
振幅 & 1.00 & 0.70 & 0.50 & 0.30 & 0.20 & 0.10 & 0.05 \\
\hline
\end{tabular}
\end{table}

\begin{table}[htbp]
\centering
\caption{金管ADSR初期値}
\begin{tabular}{llll}
\hline
項目 & 値 & 意味 & 調整方針 \\
\hline
Attack & 20 ms & 息の入りの立ち上がり & 遅すぎる場合は10 msへ短縮 \\
\hline
Decay & 40 ms & 最大音量から安定音量へ移る時間 & 音が硬い場合に延長 \\
\hline
Sustain & 最大音量の0.75倍 & 音符中の持続音量 & velocityに応じて変動 \\
\hline
Release & 60 ms & 息を止めた後の余韻 & 音の切れが悪い場合は40 msへ短縮 \\
\hline
\end{tabular}
\end{table}

\subsubsection{ドラム音色}

塩澤案ではnode\_05がドラムであるため，4パートすべてを金管で鳴らす初版設計は採用しない．
ドラムは，ノイズ成分と短い減衰を用いる．最小実装ではMIDIノート番号ごとに
キック，スネア，ハイハット風の3種類を割り当てる．

\begin{longtable}{p{0.18\linewidth}p{0.24\linewidth}p{0.46\linewidth}}
\hline
noteNumber & 音色 & 合成方針 \\
\hline
\endhead
36 & Kick & 低周波のサイン波を50--80 Hzから下げ，Release 120 msで減衰する． \\
\hline
38 & Snare & ホワイトノイズ + 180 Hz付近の短い共鳴，Release 90 ms． \\
\hline
42 & Hi-hat & 高域ノイズを短く鳴らし，Release 40 ms． \\
\hline
その他 & Click & 短いノイズで代替し，未知のnoteNumberでも無音にならないようにする． \\
\hline
\end{longtable}

\subsection{FFT解析の設計}

WBS 127は任意タスクだが，ADR-0007の「倍音データを元に根拠を持った音色作り」と対応するため，
音色パラメータを外部データで更新可能な形にしておく．

\begin{enumerate}
\item フリー素材または録音した金管単音を\texttt{tools/}配下に置く．
\item PythonでRMS包絡線とFFTを計算し，基音を1.0として第2--第7倍音を正規化する．
\item 解析結果をCSVとして保存する．列は\texttt{harmonic, amplitude}とする．
\item Processingの\texttt{ToneFactory}は，初期値をコード内に持ちつつ，CSVがあれば読み込む．
\item 音色評価では，初期値版とFFT反映版を聞き比べ，採用値を記録する．
\end{enumerate}

\section{詳細設計：処理フローの設計}

\subsection{起動処理}

\begin{enumerate}
\item \texttt{setup()}で画面，Minim，AudioOutput，UI状態，Loggerを初期化する．
\item 設定ファイルまたはキー入力から期待partIdを決定する．Mac 1--4では0x02--0x05に固定する．
\item \texttt{Serial.list()}でポート一覧を取得し，前回ポートが存在すれば\texttt{Serial}を開く．
\item 接続成功ならReadyへ，失敗ならPortSelectへ遷移する．
\item テスト音を鳴らし，スピーカと音量を確認できるようにする．
\end{enumerate}

\subsection{NOTE受信フロー}

\begin{lstlisting}[style=textdiagram, caption={Serial受信から発音までの処理フロー}]
serialEvent()
  while serial has bytes:
    reader.pushByte(serial.read())
    while reader.tryReadPacket(packet):
      if packet.version != 1:
        log version error; continue
      if packet.type != 3:
        log non-NOTE packet; continue
      if packet.gate is not 0 or 1:
        log gate error; continue
      if expectedPartId != ALL and packet.partId != expectedPartId:
        log wrong part; show warning; continue
      if packet.seq is duplicate:
        log duplicate; continue
      PartManager.handleNote(packet)
\end{lstlisting}

\subsection{ボイス管理フロー}

\begin{enumerate}
\item \texttt{gate=1}かつ\texttt{noteNumber>0}ならNoteOnとして扱う．
\item partIdから音色種別を決める．0x02--0x04は\texttt{BrassVoice}，0x05は\texttt{DrumVoice}である．
\item velocityを0.0--1.0へ正規化し，総振幅上限0.8を超えないよう\texttt{maxAmp}を決める．
\item 同一partIdの発音数が上限を超える場合は，最も古いVoiceをReleaseへ移す．
\item \texttt{durationMs}から自動NoteOff時刻を計算する．金管は最低80 ms，ドラムは最低30 msとする．
\item \texttt{gate=0}または自動NoteOff時刻到達時に，該当partIdの同一noteNumberをReleaseへ移す．
\item Release終了後，Voiceを配列から削除し，Minimの接続を解除する．
\end{enumerate}

\subsection{例外処理}

\begin{longtable}{p{0.25\linewidth}p{0.35\linewidth}p{0.28\linewidth}}
\hline
条件 & 処理 & 目的 \\
\hline
\endhead
magic不一致 & 1Bずつ破棄して次の\texttt{0x52 0x4F}を探す & バイトずれから復帰する． \\
\hline
version不一致 & パケット全体を破棄 & 旧仕様・別仕様との混在を防ぐ． \\
\hline
type不一致 & パケット全体を破棄 & CTRL/BEATを誤って発音しない． \\
\hline
seq重複 & 発音せずログのみ & 二重発音を防ぐ． \\
\hline
seq欠落 & 発音は継続し，欠落数をログ & USB Serialではまれだが検証可能にする． \\
\hline
partId不一致 & 発音せず画面警告 & MacとArduinoの接続ミスを即発見する． \\
\hline
Serial切断 & 全VoiceをReleaseしErrorへ遷移 & 無限発音を防ぐ． \\
\hline
音量過大 & 全Voiceの振幅をスケールダウン & クリップと耳に痛い音を防ぐ． \\
\hline
\end{longtable}

\subsection{主要コード断片}

\begin{lstlisting}[style=code, caption={NOTE復号の最小実装イメージ}]
class NotePacket {
  int version, type, partId, noteNumber, velocity, gate, durationMs;
  long seq, timestampMs;
}

int u16le(byte[] b, int off) {
  return (b[off] & 0xff) | ((b[off + 1] & 0xff) << 8);
}

long u32le(byte[] b, int off) {
  return ((long)b[off] & 0xff)
      | (((long)b[off + 1] & 0xff) << 8)
      | (((long)b[off + 2] & 0xff) << 16)
      | (((long)b[off + 3] & 0xff) << 24);
}

NotePacket decodeNote(byte[] frame) {
  if (u16le(frame, 0) != 0x4F52) return null;
  NotePacket p = new NotePacket();
  p.version = frame[2] & 0xff;
  p.type = frame[3] & 0xff;
  p.seq = u32le(frame, 4);
  p.timestampMs = u32le(frame, 8);
  p.partId = frame[12] & 0xff;
  p.noteNumber = frame[13] & 0xff;
  p.velocity = frame[14] & 0xff;
  p.gate = frame[15] & 0xff;
  p.durationMs = u16le(frame, 16);
  return p;
}
\end{lstlisting}

\begin{lstlisting}[style=code, caption={金管波形生成の実装イメージ}]
Waveform makeBrassWaveform() {
  return WavetableGenerator.gen10(
    4096,
    new float[] {
      1.00f, 0.70f, 0.50f, 0.30f,
      0.20f, 0.10f, 0.05f
    }
  );
}
\end{lstlisting}

\section{検証計画}

本書の仕様がルーブリック上の「実装に必要な設計」になっているかを確認するため，
Processing単体試験，Arduinoとの結合試験，音色評価を行う．

\begin{longtable}{p{0.11\linewidth}p{0.25\linewidth}p{0.34\linewidth}p{0.18\linewidth}}
\hline
ID & 試験項目 & 手順 & 合否基準 \\
\hline
\endhead
PT-1 & NOTE復号単体 & 20Bの正常フレームを固定配列で与え，\texttt{NotePacket}へ復号する． & 全フィールド一致 \\
\hline
PT-2 & フレーム同期復帰 & 先頭に破損バイトを混ぜた後，正常フレームを与える． & 次のmagicから復帰 \\
\hline
PT-3 & partId誤接続 & Mac 1設定でpartId 0x03を入力する． & 発音せず警告表示 \\
\hline
PT-4 & NoteOn/Off & gate=1後，gate=0またはdurationMs経過を入力する． & 無限発音しない \\
\hline
PT-5 & 同時発音上限 & 金管で5音以上を連続入力する． & 古いVoiceをReleaseし，上限4を維持 \\
\hline
PT-6 & ドラム音色 & partId 0x05でnote 36/38/42を入力する． & 3種類のリズム音として聞き分け可能 \\
\hline
PT-7 & 受信遅延 & 受信時刻からVoice生成までを\texttt{millis()}で100回記録する． & 最大5 ms以下 \\
\hline
PT-8 & 結合演奏 & node\_02--05とMac 4台を接続し，カエルの歌等を演奏する． & 4パートが欠落なく鳴る \\
\hline
\end{longtable}

\subsection{ルーブリック対応表}

\begin{longtable}{p{0.28\linewidth}p{0.16\linewidth}p{0.46\linewidth}}
\hline
評価項目 & 目標評価 & 本書での対応 \\
\hline
\endhead
基本設計：機能一覧とシステム構成図 & A & 機能一覧表，Mac 4台構成図，partId対応表，MOPを記載した． \\
\hline
基本設計：操作インターフェースと状態遷移 & A & キー操作，画面表示，BootからErrorまでの状態遷移と状態定義を記載した． \\
\hline
詳細設計：ハードウェア設計 & A & Arduino，Mac，USB，スピーカの数量・接続関係・動作・故障時挙動を記載した． \\
\hline
詳細設計：ソフトウェア設計 & A & ファイル構成，NOTEパケット，クラス設計，音色パラメータ，コード断片を記載した． \\
\hline
詳細設計：処理フローの設計 & A & 起動，受信，復号，発音，例外処理，ボイス停止，検証フローを一貫して記載した． \\
\hline
\end{longtable}

\section{今後の実装順序}

\begin{enumerate}
\item \texttt{pc\_app/orchestra\_processing/}に最小スケッチを作り，Minimでテスト音を鳴らす．
\item \texttt{SerialFrameReader}と\texttt{decodeNote()}を実装し，固定バイト列でPT-1/2を通す．
\item \texttt{PartManager}を実装し，partId検査，NoteOn/Off，自動停止を通す．
\item 金管3パートの\texttt{BrassVoice}を実装し，既存の\texttt{HackInstrument}をADSR対応へ拡張する．
\item node\_05用の\texttt{DrumVoice}を実装する．
\item UIとCSVログを実装し，Mac 4台の接続ミスが見える状態にする．
\item Arduino側の\texttt{NoteSenderModule}と結合し，PT-8を実施する．
\end{enumerate}

\section{生成AIの利用について}

本改訂案の作成にあたり，生成AIを利用した．
利用範囲は，既存資料のレビュー結果を踏まえた章立ての整理，
ルーブリック項目との対応確認，および文章案の作成である．
ただし，採用する前提条件，NOTEパケット仕様，Mac 4台構成，partId対応，
試験項目については，リポジトリ内資料およびチーム内レビューを確認したうえで著者が判断する．

\end{document}
